\section{Conclusion}
\label{sec:conclusion}

This work presented an end-to-end analysis of AI-assisted academic research across the complete lifecycle. We organized the field into four epistemological phases: \emph{Creation}, \emph{Writing}, \emph{Validation}, and \emph{Dissemination}, and eight stages spanning ideation, literature review, coding and experiments, tables and figures, paper writing, peer review, rebuttal and revision, and Paper2X dissemination. This lifecycle framing connects tools that are often studied in isolation and exposes where current systems succeed, where they fail, and how errors propagate across stage boundaries.

The central finding is that AI systems are increasingly capable of producing research artifacts, but remain less reliable at verifying their scientific meaning. Across the lifecycle, plausible outputs can conceal deeper failures: ideas may weaken after execution, retrieved evidence may be misrepresented, executable code may implement the wrong algorithm, fluent manuscripts may lack argumentative depth, reviews may miss methodological flaws, rebuttals may promise unfulfilled revisions, and dissemination artifacts may overstate claims. The core bottleneck is therefore not generation alone, but maintaining novelty, faithfulness, reproducibility, and accountability across the research process.

The most credible path forward is human-governed AI-assisted research. AI should reduce mechanical friction in retrieval, drafting, coding, visualization, review support, and dissemination, while researchers retain ownership over judgment, interpretation, experimental design, argumentation, and final responsibility. Future systems should maintain provenance across artifacts, use retrieval and execution grounding wherever possible, support human checkpoints at phase boundaries, and make AI involvement transparent. If developed with these principles, AI can amplify human creativity and rigor; without them, it risks scaling the production of plausible but unreliable research artifacts.


\subsection*{Acknowledgments}
We thank Josh Susskind for insightful discussions and careful proofreading of this manuscript. 

We also thank the researchers and open-source contributors whose systems, benchmarks, datasets, and technical reports made this survey possible, as well as the broader community for ongoing discussions on responsible AI use, research integrity, peer review, and the future of scientific work.


\subsection*{Responsible Use and Limitations}
This work is intended to inform responsible use of AI-assisted research tools, not to endorse replacing human scientific judgment with full automation. Current systems are most reliable when used to assist retrieval, drafting, coding, visualization, review support, and dissemination, while humans retain responsibility for novelty, interpretation, verification, authorship, and accountability. Because the field evolves rapidly, this paper should be read as a structured snapshot through our search cutoff, and AI-generated research outputs should be independently verified before scholarly use.